%% ----------------------------------------------------------------------------
%%  §11.4  La serie geometrica
%% ----------------------------------------------------------------------------
\section{La serie geometrica}
\label{sec:serie-geometrica}

Per ogni numero reale $x$ consideriamo la \emph{serie geometrica} di \emph{ragione} $x$
\begin{equation}\label{eq:serie-geometrica-def}
  \sum_{k=0}^{\infty}x^{k}=1+x+x^{2}+\dots+x^{n}+\cdots .
\end{equation}
È la serie più importante di tutta la teoria: è una delle pochissime di cui si conosce
esplicitamente la somma parziale, e servirà da termine di paragone per molte altre.

\begin{teorema}[Serie geometrica]\label{teo:serie-geometrica}\index{serie!geometrica}
La serie geometrica di ragione $x$ è convergente se $\abs x<1$, divergente
positivamente se $x\ge1$, indeterminata se $x\le-1$. Inoltre
\begin{equation}\label{eq:serie-geometrica}
  \sum_{k=0}^{\infty}x^{k}=\frac1{1-x} \qquad\text{se } -1<x<1 .
\end{equation}
\end{teorema}

\begin{schema}
\Ipotesi{$x\in\R$.}
\Tesi{$\sum x^{k}=\frac1{1-x}$ se $\abs x<1$; \ $\sum x^{k}=+\infty$ se $x\ge1$; \ $\sum x^{k}$
indeterminata se $x\le-1$.}
\end{schema}

\begin{strategia}
Per $x\neq1$ la somma parziale è nota esplicitamente, $s_{n}=\frac{1-x^{n+1}}{1-x}$, e il suo
comportamento dipende soltanto da quello di $x^{n+1}$, che conosciamo dal capitolo sulle
successioni. Il caso $x=1$ si tratta a parte.
\end{strategia}

\begin{dimostrazione}
Se $x=1$, $s_{n}=n+1\to+\infty$. Se $x\neq1$, per la \cref{prop:geometrica}
\begin{equation}\label{eq:somma-parziale-geometrica}
  s_{n}=\sum_{k=0}^{n}x^{k}=\frac{1-x^{n+1}}{1-x} .
\end{equation}
Per il \cref{teo:lim-geometrica}, $x^{n+1}\to0$ se $-1<x<1$, $x^{n+1}\to+\infty$ se $x>1$, e
$x^{n+1}$ non ha limite se $x\le-1$. Nel primo caso $s_{n}\to\frac1{1-x}$. Nel secondo, poiché
$1-x<0$, $s_{n}=\frac{x^{n+1}-1}{x-1}\to+\infty$. Nel terzo, se $s_{n}$ avesse limite $\ell$, anche
$x^{n+1}=1-(1-x)s_{n}$ avrebbe limite $1-(1-x)\ell$ (finito o infinito), il che è falso.
\end{dimostrazione}

Moltiplicando per una costante (\cref{prop:linearita-serie}) e trascurando i primi
termini, si ottiene la forma in cui la serie geometrica si usa più spesso: per $\abs x<1$,
$a\in\R$ e $m\ge0$ intero,
\begin{equation}\label{eq:geometrica-generale}
  \sum_{k=m}^{\infty}a\,x^{k}=a\,x^{m}+a\,x^{m+1}+\dots=\frac{a\,x^{m}}{1-x} ,
\end{equation}
cioè «il primo termine diviso per uno meno la ragione». Ad esempio
$\frac12+\frac14+\frac18+\cdots=\frac{1/2}{1-1/2}=1$, come mostra la \cref{fig:quadrato-geometrico}.

\begin{figure}[htb]
\centering
\begin{tikzpicture}[scale=3.2]
  \fill[blu!35] (0,0) rectangle (0.5,1);
  \fill[blu!25] (0.5,0) rectangle (1,0.5);
  \fill[blu!17] (0.5,0.5) rectangle (0.75,1);
  \fill[blu!11] (0.75,0.5) rectangle (1,0.75);
  \fill[blu!7]  (0.75,0.75) rectangle (0.875,1);
  \fill[blu!4]  (0.875,0.75) rectangle (1,0.875);
  \draw[blu] (0,0) rectangle (1,1);
  \draw[blu] (0.5,0) -- (0.5,1) (0.5,0.5) -- (1,0.5) (0.75,0.5) -- (0.75,1)
    (0.75,0.75) -- (1,0.75) (0.875,0.75) -- (0.875,1) (0.875,0.875) -- (1,0.875);
  \node[font=\small] at (0.25,0.5) {$\frac12$};
  \node[font=\small] at (0.75,0.25) {$\frac14$};
  \node[font=\footnotesize] at (0.625,0.75) {$\frac18$};
  \node[font=\scriptsize] at (0.875,0.625) {$\frac1{16}$};
  \node[font=\tiny] at (0.8125,0.875) {$\frac1{32}$};
\end{tikzpicture}
\caption{La serie geometrica $\frac12+\frac14+\frac18+\cdots=1$: ogni rettangolo è metà della
parte del quadrato unitario non ancora coperta, e i rettangoli riempiono l'intero
quadrato. Le somme parziali sono $1-\frac1{2^{n}}$.}
\label{fig:quadrato-geometrico}
\end{figure}

\subsection{I numeri decimali periodici}

Con la serie geometrica possiamo giustificare la regola, nota dalla scuola media, con
cui un numero decimale periodico si trasforma nella sua \emph{frazione generatrice}. Un
allineamento decimale infinito $0{,}c_{1}c_{2}c_{3}\ldots$, con cifre $c_{k}\in\{0,1,\dots,9\}$, è per
definizione la somma della serie a termini non negativi $\sum_{k=1}^{\infty}c_{k}10^{-k}$, che
converge perché le sue somme parziali sono limitate da $\sum_{k=1}^{\infty}9\cdot10^{-k}=1$.

\begin{proposizione}[Frazione generatrice]\label{prop:periodici}
Sia $0{,}\overline{c_{1}c_{2}\ldots c_{p}}$ un numero decimale periodico semplice con periodo di $p$
cifre, e sia $C=c_{1}c_{2}\ldots c_{p}$ l'intero che ha quelle cifre. Allora
\begin{equation}\label{eq:periodici}
  0{,}\overline{c_{1}c_{2}\ldots c_{p}}=\frac C{10^{p}-1}=\frac C{99\ldots9} ,
\end{equation}
dove il denominatore $99\ldots9$ è formato da $p$ cifre uguali a $9$. In particolare, ogni
numero decimale periodico è razionale.
\end{proposizione}

\begin{schema}
\Ipotesi{$0{,}\overline{c_{1}\ldots c_{p}}=\sum_{k=1}^{\infty}c_{k}10^{-k}$ con $c_{k+p}=c_{k}$; \ $C=c_{1}c_{2}\ldots c_{p}$.}
\Tesi{$0{,}\overline{c_{1}\ldots c_{p}}=\frac C{10^{p}-1}$.}
\end{schema}

\begin{strategia}
Le cifre si ripetono a blocchi di $p$: ogni blocco vale $C$ volte una potenza di
$10^{-p}$, e raggruppando i termini a blocchi la serie delle cifre diventa una serie
geometrica di ragione $10^{-p}$, di cui conosciamo la somma.
\end{strategia}

\begin{dimostrazione}
Il numero è la somma della serie in cui il blocco di cifre si ripete: il primo blocco
vale $C\cdot10^{-p}$, il secondo $C\cdot10^{-2p}$, e così via. Raggruppando le cifre a blocchi
di $p$ (\cref{oss:parentesi}) e usando la \eqref{eq:geometrica-generale} con $x=10^{-p}$,
\[
  0{,}\overline{c_{1}\ldots c_{p}}=\sum_{j=1}^{\infty}C\,10^{-jp}=\frac{C\,10^{-p}}{1-10^{-p}}=\frac C{10^{p}-1} .
\]
Un numero periodico misto si riconduce a questo caso moltiplicando per una potenza di
$10$ e sottraendo l'antiperiodo, che è un numero razionale.
\end{dimostrazione}

Ad esempio $0{,}\overline{13}=\frac{13}{99}$, e $0{,}\overline{9}=\frac99=1$: l'allineamento $0{,}999\ldots$
\emph{non} è «un po' meno» di $1$, ma è esattamente $1$, perché la sua somma è il
\emph{limite} delle somme parziali $1-10^{-n}$. Molte perplessità sull'uguaglianza
$0{,}\overline9=1$ nascono proprio dal confondere una serie con le sue somme parziali.

\subsection{Il paradosso di Achille}

\begin{idea}[Achille e la tartaruga, in volo]
Un velivolo $A$ insegue alla velocità $v_{A}$ un velivolo $B$ che lo precede di una
distanza $d$ e vola sulla stessa rotta alla velocità $v_{B}<v_{A}$. Ragioniamo come
Zenone (\cref{fig:zenone}). Nel tempo $t_{0}=\frac d{v_{A}}$ che $A$ impiega a raggiungere la
posizione iniziale di $B$, questo avanza di $d\,\frac{v_{B}}{v_{A}}$; nel tempo $t_{1}=t_{0}\frac{v_{B}}{v_{A}}$ che
$A$ impiega a colmare questo nuovo distacco, $B$ avanza di $d\bigl(\frac{v_{B}}{v_{A}}\bigr)^{2}$; e così
via. I tempi $t_{k}=t_{0}q^{k}$, con $q=\frac{v_{B}}{v_{A}}<1$, sono infiniti, ma la loro somma è la
serie geometrica
\[
  T=\sum_{k=0}^{\infty}\frac d{v_{A}}\,q^{k}=\frac d{v_{A}}\cdot\frac1{1-\frac{v_{B}}{v_{A}}}=\frac d{v_{A}-v_{B}} ,
\]
che è finita, ed è proprio il tempo che si ottiene ragionando con la velocità relativa
$v_{A}-v_{B}$. Ad esempio, con $d=20$~km, $v_{A}=900$~km/h e $v_{B}=750$~km/h, il
raggiungimento avviene dopo $T=\frac{20}{150}$~h, cioè $8$~minuti. Dal punto di vista matematico, il
paradosso nasce dal supporre che infiniti intervalli di tempo debbano avere una durata
complessiva infinita. (Aristotele, che riferisce l'argomento, rispondeva distinguendo
l'infinito «in atto» da quello «in potenza»: un dibattito filosofico che la definizione
di somma di una serie non esaurisce, ma rende innocuo per il calcolo.)
\end{idea}

\begin{figure}[htb]
\centering
\begin{tikzpicture}[x=0.95cm]
  \draw[grigio, -{Stealth[length=2mm]}] (-0.3,0) -- (10.6,0) node[above, font=\scriptsize] {rotta};
  \foreach \x/\lab in {0/{$A_{0}$}, 4/{$A_{1}=B_{0}$}, 6.4/{$A_{2}=B_{1}$}, 7.84/{$B_{2}$}}
    \fill (\x,0) circle (1.4pt) node[below=2pt, font=\scriptsize] {\lab};
  \fill (8.704,0) circle (1pt);
  \fill (9.222,0) circle (1pt);
  \fill (9.533,0) circle (0.8pt);
  \draw[rame, thick] (10,-0.12) -- (10,0.12);
  \node[rame, font=\scriptsize, below=2pt] at (10,0) {incontro};
  \draw[blu, -{Stealth[length=1.6mm]}] (0,0.35) -- node[above, font=\scriptsize] {$t_{0}$} (4,0.35);
  \draw[blu, -{Stealth[length=1.6mm]}] (4,0.35) -- node[above, font=\scriptsize] {$t_{1}=qt_{0}$} (6.4,0.35);
  \draw[blu, -{Stealth[length=1.6mm]}] (6.4,0.35) -- node[above, font=\scriptsize] {$t_{2}$} (7.84,0.35);
  \draw[blu, -{Stealth[length=1.6mm]}] (7.84,0.35) -- (8.704,0.35);
\end{tikzpicture}
\caption{Il ragionamento di Zenone per l'inseguimento ($q=\frac{v_{B}}{v_{A}}=0{,}6$ nel disegno):
$A$ raggiunge successivamente le posizioni $B_{0},B_{1},B_{2},\dots$ occupate da $B$ un passo
prima. I tratti, e i tempi, formano una progressione geometrica di ragione $q$, e la loro
somma è finita: l'incontro avviene a distanza $\frac d{1-q}$ dal punto di partenza di $A$.}
\label{fig:zenone}
\end{figure}

\begin{python}[Serie geometrica e frazioni generatrici]
from fractions import Fraction as F

def periodico(antiperiodo, periodo):
    """Frazione di 0,antiperiodo(periodo), es. ('1', '6') -> 1/6."""
    m, p = len(antiperiodo), len(periodo)
    A = int(antiperiodo) if antiperiodo else 0
    C = int(periodo)
    # 0,a(c) = (A + C/(10^p - 1)) / 10^m
    return (A + F(C, 10**p - 1)) / 10**m

for a, c in (("", "13"), ("", "9"), ("1", "6"), ("", "142857")):
    print(f"0,{a}({c}) = {periodico(a, c)}")

x = 0.6                                    # somme parziali, q = 0,6
s = [sum(x**k for k in range(n + 1)) for n in (5, 10, 20, 40)]
print("s_n:", [round(v, 8) for v in s], " 1/(1-x) =", 1 / (1 - x))
\end{python}
\begin{risultato}
0,(13) = 13/99
0,(9) = 1
0,1(6) = 1/6
0,(142857) = 1/7
s_n: [2.38336, 2.49093007, 2.49994516, 2.5]  1/(1-x) = 2.5
\end{risultato}

\begin{notastorica}[Eulero e le «somme» delle serie divergenti]
La formula \eqref{eq:serie-geometrica} vale soltanto per $-1<x<1$, ma per tutto il
Settecento fu usata spesso senza questa restrizione. Con $x=-1$ se ne ricava
$1-1+1-1+\cdots=\frac12$: il camaldolese Guido Grandi, professore a Pisa, vi vide nel 1703
un'immagine della creazione dal nulla, e Leibniz, in una lettera a Christian Wolff
pubblicata nel 1713, difese il valore $\frac12$ osservando che le somme parziali valgono $1$
e $0$ «con uguale probabilità». Con $x=2$ si ottiene addirittura $-1=1+2+4+8+\cdots$.
Eulero discusse proprio questa uguaglianza nella memoria \emph{De seriebus divergentibus}
(1760), e propose di chiamare somma di una serie l'espressione finita dal cui sviluppo
essa nasce, qui $\frac1{1-x}$: nelle sue mani la regola condusse a scoperte profonde, in mani
meno esperte a contraddizioni. Cauchy, nel \emph{Cours d'analyse} (1821), tagliò il nodo
con la definizione che abbiamo adottato: una serie divergente non ha somma, e le
operazioni sulle serie sono lecite soltanto per quelle convergenti. Nel 1826 il giovane
Niels Henrik Abel scriveva a un amico che le serie divergenti sono «qualcosa di funesto»,
su cui è una vergogna fondare una dimostrazione. Più tardi si capì come recuperare in
modo coerente almeno una parte di quei valori: Ernesto Cesàro (1890) propose di
sostituire alle somme parziali le loro medie aritmetiche, che per la serie di Grandi
tendono proprio a $\frac12$. Nacque così la teoria della \emph{sommabilità}, che non
contraddice la definizione di Cauchy, ma la estende.
\end{notastorica}
